\section{วิธีดำเนินการวิจัย}

การดำเนินงานวิจัยนี้แบ่งออกเป็น 5 ส่วนหลัก เพื่อให้การพัฒนาแอปพลิเคชันเป็นไปอย่างมีระบบและครอบคลุมทั้งด้านเทคโนโลยี การจัดการข้อมูล และการประเมินผลลัพธ์

\subsection{การศึกษาวรรณกรรมและทฤษฎีที่เกี่ยวข้อง}

ขั้นตอนนี้มีวัตถุประสงค์เพื่อรวบรวมและวิเคราะห์แนวคิด เทคนิค และเครื่องมือที่ทันสมัยในการพัฒนาระบบแนะนำการท่องเที่ยว (Tourism Recommender Systems - TRS) และระบบอัจฉริยะ

\subsubsection{การศึกษาระบบแนะนำการท่องเที่ยวและการวางแผนการเดินทาง}

เน้นการศึกษาแนวคิดพื้นฐานของระบบ TRS และวิวัฒนาการไปสู่ระบบที่พิจารณาหลายภาคส่วน (Multistakeholder Recommendation) โดยเฉพาะอย่างยิ่งการบูรณาการด้านความยั่งยืน (Sustainability) เข้าไปในการวางแผนการเดินทาง ซึ่งความยั่งยืนดังกล่าวรวมถึงผลกระทบทางเศรษฐกิจ สังคม และสิ่งแวดล้อม แนวคิดหลักคือการพัฒนาระบบให้สามารถเพิ่มประสิทธิภาพของการเดินทาง โดยพิจารณาปัจจัยหลายด้านพร้อมกัน (Multi-Objective Optimization) ได้แก่ ต้นทุน เวลา ความชอบของผู้ใช้งาน และความยั่งยืนด้านสิ่งแวดล้อม

\subsubsection{การศึกษาเทคโนโลยีปัญญาประดิษฐ์สำหรับการประมวลผลภาษาธรรมชาติ}

ศึกษานวัตกรรมที่ใช้ LLMs และการประมวลผลภาษาธรรมชาติ (Natural Language Processing - NLP) สำหรับการสร้างคำแนะนำและการตอบโต้กับผู้ใช้งาน โดยเน้นที่สถาปัตยกรรม RAG ซึ่งเป็นแนวทางที่มีประสิทธิภาพในการจัดการข้อมูลการท่องเที่ยวที่มีการเปลี่ยนแปลงอยู่เสมอ (Dynamic Data) โดยไม่จำเป็นต้องปรับจูนโมเดล LLM บ่อยครั้ง

\subsubsection{การศึกษาสถาปัตยกรรมระบบ}

ศึกษาสถาปัตยกรรมไมโครเซอร์วิส (Microservices Architecture) ซึ่งช่วยให้ระบบมีความยืดหยุ่น ปรับขนาดได้ดี (Scalability) และมีความทนทานต่อความผิดพลาด (Fault Tolerance) ซึ่งสถาปัตยกรรมนี้เหมาะสำหรับแอปพลิเคชันปัญญาประดิษฐ์ที่มีการประมวลผลข้อมูลที่ซับซ้อนและมีการเปลี่ยนแปลงแบบเรียลไทม์

\subsection{การเก็บรวบรวมและเตรียมข้อมูล}

การเตรียมข้อมูลถือเป็นรากฐานสำคัญในการสร้างระบบอัจฉริยะที่แม่นยำและน่าเชื่อถือ

\subsubsection{การรวบรวมข้อมูลแหล่งท่องเที่ยวและกิจกรรม}

ข้อมูลจะถูกรวบรวมจาก Google API เป็นหลัก ครอบคลุมข้อมูลสำคัญของสถานที่ท่องเที่ยว โรงแรม ร้านอาหาร โดยข้อมูลรวมถึงลักษณะเฉพาะของสถานที่

\subsubsection{การเตรียมและจัดเก็บข้อมูล}

\begin{enumerate}[leftmargin=1.8cm]
\item การทำความสะอาดและจัดรูปแบบข้อมูล (Data Preprocessing): ข้อมูลดิบที่รวบรวมมาจะถูกนำมาทำความสะอาด จัดรูปแบบ และทำให้เป็นมาตรฐาน (Normalizing) เพื่อให้มั่นใจในคุณภาพและความถูกต้องก่อนนำไปใช้ในโมเดลปัญญาประดิษฐ์
\item การแบ่งส่วนข้อมูลและการสร้างเวกเตอร์ (Chunking and Vector Embeddings): สำหรับการใช้งานร่วมกับ RAG ข้อมูลจะถูกแบ่งออกเป็นส่วนย่อยที่มีความสอดคล้องทางตรรกะ (Logically Coherent Pieces) จากนั้นจะทำการสร้างเวกเตอร์ฝังตัว (Vector Embeddings) โดยใช้โมเดลที่เหมาะสม
\item การจัดเก็บในฐานข้อมูลเวกเตอร์: เวกเตอร์ที่สร้างขึ้นจะถูกจัดเก็บไว้ในฐานข้อมูลเวกเตอร์ (Vector Database) เพื่อใช้ในการคำนวณความคล้ายคลึงกัน (Similarity) ในขั้นตอนการดึงข้อมูล
\end{enumerate}

\subsection{การออกแบบและพัฒนาระบบ}

การพัฒนาระบบจะใช้สถาปัตยกรรมไมโครเซอร์วิส เพื่อแยกส่วนการทำงานหลัก ๆ ออกจากกัน โดยเน้นการทำงานของส่วน Trip Planning (Itinerary Service) และแชตบอต (RAG/AIML)

\subsubsection{การออกแบบสถาปัตยกรรมไมโครเซอร์วิส}

การพัฒนาระบบในงานวิจัยนี้ใช้สถาปัตยกรรมแบบ Microservices เพื่อให้แต่ละฟังก์ชันของระบบสามารถทำงานได้อย่างเป็นอิสระต่อกัน ส่งผลให้ระบบมีความยืดหยุ่นและสามารถขยายตัวได้ง่าย นอกจากนี้ยังมีการใช้ API Gateway เป็นจุดศูนย์กลางในการรับคำขอจากผู้ใช้งาน (Client) และทำหน้าที่จัดการการส่งต่อคำขอไปยังบริการต่าง ๆ ภายในระบบ

ในส่วนของการสื่อสารระหว่างบริการ (Service-to-Service Communication) ระบบใช้ RESTful APIs สำหรับการทำงานแบบ Synchronous Communication ซึ่งต้องการการตอบสนองทันที และใช้ Message Broker เช่น RabbitMQ หรือ Apache Kafka สำหรับการสื่อสารแบบ Asynchronous Communication เพื่อเพิ่มประสิทธิภาพและลดการพึ่งพากันโดยตรงระหว่างบริการ

\subsubsection{การพัฒนาโมเดลวางแผนการเดินทาง (Trip Planning Optimization)}

Itinerary Service จะใช้ LLM ร่วมกับ RAG เป็นกลไกหลักในการสร้างแผนการเดินทางที่เหมาะสม โดยพิจารณาปัจจัยต่างๆ ดังนี้

ขั้นตอนการทำงานของ LLM + RAG มีดังนี้
\begin{enumerate}[leftmargin=1.8cm]
\item Query Understanding: รับคำสั่งจากผู้ใช้ ได้แก่ ระยะเวลา งบประมาณ และความสนใจของผู้ใช้ต่างๆ
\item Retrieval (RAG): สร้าง Query Embedding จากข้อมูลที่ผู้ใช้ระบุ แล้วค้นหาความคล้ายคลึง (Cosine Similarity) กับข้อมูลสถานที่ท่องเที่ยว ร้านอาหาร และที่พัก ที่เก็บไว้ใน Vector Database
\item Prompt Construction \& Generation: รวม User Query และข้อมูลที่ดึงมาได้เป็น Prompt แล้วส่งให้ LLM สังเคราะห์เป็นแผนการเดินทางฉบับสมบูรณ์พร้อมคำอธิบายที่เข้าใจง่าย
\end{enumerate}

\subsubsection{การพัฒนาระบบแชตบอตและคำแนะนำ (Recommendation System)}

\begin{enumerate}[leftmargin=1.8cm]
\item \textbf{Information Retrieval:} ระบบจะรับคำสั่งที่เป็นภาษาธรรมชาติ (Natural Language Query) จากผู้ใช้ และเปลี่ยนเป็น Query Embedding (paraphrase-multilingual-MiniLM-L12-v2) เพื่อค้นหาความคล้ายคลึงกัน (Cosine Similarity) กับเอกสารที่เก็บไว้ใน Vector Store ที่สร้างขึ้นจะสอดคล้องกับวัตถุประสงค์ด้านความยั่งยืน
\item \textbf{Response Generation:} ใช้ LLMs (gemini-2.5-flash) เพื่อสร้างคำแนะนำที่เป็นข้อความ (Textual Information) พร้อมคำอธิบาย (Explanation) ที่สมเหตุสมผล โดยสร้างตามปัจจัย
\end{enumerate}

\subsubsection{การออกแบบและพัฒนาเว็บแอปพลิเคชัน (Web Application Design and Development)}

เพื่อให้ผู้ใช้งานสามารถเข้าถึงระบบการแนะนำการท่องเที่ยวได้อย่างสะดวก จะมีการออกแบบและพัฒนาเว็บแอปพลิเคชันที่เน้น ความเป็นมิตรต่อผู้ใช้ (User-Friendly) และ การตอบสนองที่รวดเร็ว (Responsive Design) โดยมีขั้นตอนดังนี้

\begin{enumerate}[leftmargin=1.8cm]
\item \textbf{การออกแบบ UX/UI}
  \begin{enumerate}[leftmargin=1.8cm]
  \item กำหนดโฟลว์การใช้งาน (User Flow) ตั้งแต่การเข้าสู่ระบบ การป้อนความต้องการท่องเที่ยว ไปจนถึงการรับแผนการเดินทางที่ระบบสร้างให้
  \item ออกแบบ Wireframe และ Mockup ของหน้าจอสำคัญ เช่น หน้า Home, หน้า Trip Planning, หน้า Recommendation Results, หน้า Sustainability Report
  \item ใช้หลักการ User-Centered Design (UCD) เพื่อให้สอดคล้องกับพฤติกรรมของนักท่องเที่ยว
  \end{enumerate}
\item \textbf{การพัฒนา Frontend}
  \begin{enumerate}[leftmargin=1.8cm]
  \item ใช้เฟรมเวิร์ก Node.JS เพื่อสร้างอินเทอร์เฟซที่โต้ตอบได้ (Interactive UI)
  \item รองรับการใช้งานข้ามอุปกรณ์ (Responsive Web Design) ทั้งบน Desktop, Tablet และ Mobile
  \end{enumerate}
\item \textbf{การพัฒนา Backend และการเชื่อมต่อกับไมโครเซอร์วิส}
  \begin{enumerate}[leftmargin=1.8cm]
  \item ออกแบบ API สำหรับดึงข้อมูลจาก Itinerary Service, Recommendation System และ Sustainability Service
  \item จัดการ Authentication \& Authorization เพื่อรักษาความปลอดภัยของผู้ใช้
  \item พัฒนา Dashboard สำหรับผู้ดูแลระบบ (Admin) เพื่อจัดการข้อมูลแหล่งท่องเที่ยวและติดตามการใช้งานระบบ
  \end{enumerate}
\item \textbf{การทดสอบด้านการใช้งาน (Usability Testing)}
  \begin{enumerate}[leftmargin=1.8cm]
  \item ทดสอบกับกลุ่มผู้ใช้เป้าหมาย (นักศึกษา/นักท่องเที่ยวในจังหวัดขอนแก่น)
  \item ใช้เกณฑ์ประเมินด้านความง่ายในการใช้งาน (Ease of Use), ความพึงพอใจ (Satisfaction) และความเข้าใจในการนำเสนอข้อมูล
  \end{enumerate}
\end{enumerate}

\subsection{การทดสอบและประเมินผลระบบ}

\subsubsection{การทดสอบประสิทธิภาพของระบบวางแผนการเดินทาง}

จะมีการวัดผลลัพธ์หลักหลายด้าน (Performance Metrics)
\begin{enumerate}[leftmargin=1.8cm]
\item เวลาตอบสนอง (Response Time): วัดระยะเวลาที่ระบบใช้ในการสร้างแผนการเดินทางที่เหมาะสมที่สุด (Optimized Itinerary)
\item ความแม่นยำในการตอบสนองความต้องการของผู้ใช้ (Accuracy in Meeting User Preferences): วัดสัดส่วนของความต้องการที่ผู้ใช้ระบุ (เช่น งบประมาณ ประเภทโรงแรม) ที่ได้รับการบรรจุไว้ในแผนการเดินทาง
\end{enumerate}

\subsubsection{การประเมินผลโมเดล RAG/แชตบอต}

\begin{enumerate}[leftmargin=1.8cm]
\item ความเกี่ยวข้องของคำตอบ (Answer Relevance): ประเมินว่าคำตอบที่สร้างโดย LLM นั้นตอบคำถามของผู้ใช้ได้ดีเพียงใด โดยอาจใช้ LLM อื่น ๆ (เช่น GPT-4o-mini, Claude-3-5-sonnet) เป็นผู้ประเมิน (Judge) ในระดับ 0--100 (LLM as a Judge)
\item ความน่าเชื่อถือ (Faithfulness): นับจำนวนคำตอบที่อยู่นอกบริบท (Out-of-Context - OC) ซึ่งบ่งชี้ถึงการเกิด Hallucination ของโมเดล
\end{enumerate}

\subsubsection{การประเมินความพึงพอใจของผู้ใช้}

ใช้แบบสอบถามเพื่อประเมินความรู้สึกและทัศนคติของผู้ใช้เมื่อมีปฏิสัมพันธ์กับระบบปัญญาประดิษฐ์ โดยเฉพาะอย่างยิ่งการวัดความพึงพอใจโดยรวม ความรู้สึกด้านบวกหรือด้านลบ และการยอมรับทางเทคโนโลยี สำหรับแชตบอต โดยประเมินความพึงพอใจของผู้ใช้ (User Satisfaction) ความแม่นยำ (Accuracy) Precision, Recall และ F-score

\subsection{การวิเคราะห์และสรุปผล}

\subsubsection{การวิเคราะห์ผลลัพธ์ของโมเดล}

เปรียบเทียบผลการประเมินที่ได้ในส่วนที่ 4 โดยเน้นการวิเคราะห์เชิงลึก
\begin{enumerate}[leftmargin=1.8cm]
\item Trip Planning: วิเคราะห์ว่าการใช้ LLM ในสถาปัตยกรรม Microservices สามารถปรับปรุงคุณภาพของแผนการเดินทาง (Itinerary Quality) ได้อย่างมีประสิทธิภาพตามเป้าหมายที่ตั้งไว้หรือไม่
\item แชตบอต: วิเคราะห์ว่าการรวมกลไก RAG ช่วยให้ระบบสามารถให้คำแนะนำที่สอดคล้องกับความยั่งยืนมากขึ้นโดยไม่ลดทอนคุณภาพและความเกี่ยวข้องของคำตอบ (Answer Relevance) หรือไม่
\end{enumerate}

\subsubsection{การวิเคราะห์เชิงพฤติกรรมและการยอมรับของผู้ใช้}

วิเคราะห์ข้อมูลความพึงพอใจของผู้ใช้ โดยพิจารณาว่าผู้ใช้รับรู้ถึงประโยชน์หลักของการใช้ปัญญาประดิษฐ์อย่างไร (เช่น การเข้าถึงข้อมูลที่รวดเร็วขึ้น กระบวนการจองที่ง่ายขึ้น เวลาการรอสั้นลง) และระบุข้อกังวลหลักของผู้ใช้ (เช่น ปัญหาความเป็นส่วนตัวและความปลอดภัยของข้อมูล การพึ่งพาเทคโนโลยีสูง) ซึ่งข้อมูลเหล่านี้สามารถนำไปสู่การปรับปรุงความโปร่งใสของอัลกอริทึมและการสร้างความเชื่อมั่นในระบบ

\subsubsection{การสรุปผลและข้อเสนอแนะ}

สรุปผลลัพธ์หลักที่แสดงให้เห็นถึงความสำเร็จของการบูรณาการปัญญาประดิษฐ์ในการวางแผนการเดินทางและการแนะนำที่ขับเคลื่อนด้วยความยั่งยืน พร้อมระบุข้อจำกัดที่พบในการวิจัย (เช่น ปัญหา Cold-start สำหรับผู้ใช้ใหม่ หรือขนาดชุดข้อมูลที่จำกัดสำหรับ AIML) และเสนอแนวทางสำหรับการวิจัยในอนาคต เช่น การขยายขอบเขตความรู้ให้ครอบคลุมตัวชี้วัดความยั่งยืนที่หลากหลายขึ้น (เช่น Carbon Footprint และผลกระทบทางเศรษฐกิจในท้องถิ่น) หรือการพัฒนาเป็นระบบแนะนำแบบสนทนาที่ซับซ้อนยิ่งขึ้น (Conversational Recommender System)
